% Compile with: uplatex file.tex && dvipdfmx file.dvi
%           or: latexmk -uplatex file.tex
%\documentclass[uplatex,dvipdfmx]{jsarticle}
\documentclass[dvipdfmx,a4paper,11pt]{jsarticle}
\usepackage{amsmath, amssymb}
\usepackage{tikz}

\title{scenario\_larmor\_precession}
\date{\today}

\begin{document}
\maketitle

磁場中のスピン1/2のラーモア歳差運動を扱う。スピン演算子 $S$（成分 $S_x,S_y,S_z$）を準備する。

S をスピン演算子として定義する

\[
  S, S^{x}, S^{y}, S^{z}
\]

磁気回転比 $\gamma$ と磁場の強さ $B$ を実数スカラーパラメータとして準備する。

$\gamma$ と B を実数スカラーとして定義する

\[
  \gamma, B
\]

Zeeman ハミルトニアン $H=-\gamma BS_z$（磁場中のスピンのエネルギー）を登録する。

$H = -\gamma B Sz_1$

\begin{equation}
  H = -\gamma B Sz_1
  \label{eq:E1}
\end{equation}

ハミルトニアン $H$ とスピンの $x$ 成分 $S_x$ の交換子を計算する。この交換子から $S_x$ の時間発展の一部が決まる（QuantumAlgebraは $[S_x,S_y]=2iS_z$ というPauli行列規約を使うため、教科書の $iS_z$ の2倍の係数になる）。

H と $Sx_1$ の交換子を計算して

\begin{equation}
  - 2 B \gamma \mathit{i} Sy_1
  \label{eq:E2}
\end{equation}

同様に、ハミルトニアン $H$ とスピンの $y$ 成分 $S_y$ の交換子を計算する。

H と $Sy_1$ の交換子を計算して

\begin{equation}
  2 B \gamma \mathit{i} Sx_1
  \label{eq:E3}
\end{equation}

ハイゼンベルク方程式 $dS_x/dt=i[H,S_x]$ により、$S_x$ の運動方程式を求める。歳差運動によって $S_y$ に比例する項が現れるはずである。

$Sx_1$ の Heisenberg 方程式 H を計算して

\begin{equation}
  \frac{d Sx_1}{d t} = 2 B \gamma Sy_1
  \label{eq:E4}
\end{equation}

同様に $S_y$ の運動方程式を求める。$S_x$ に比例する項が符号を変えて現れるはずである。

$Sy_1$ の Heisenberg 方程式 H を計算して

\begin{equation}
  \frac{d Sy_1}{d t} = - 2 B \gamma Sx_1
  \label{eq:E5}
\end{equation}

$S_z$ の運動方程式を求める。$H$ が $S_z$ のみで書かれているため、$S_z$ は保存量でありゼロになるはずである。

$Sz_1$ の Heisenberg 方程式 H を計算して

\begin{equation}
  \frac{d Sz_1}{d t} = 0
  \label{eq:E6}
\end{equation}

$H$ と $S_z$ の交換子を直接計算し、ステップ8の結果（保存量であること）を別の方法で確認する。

H と $Sz_1$ の交換子を計算して

\begin{equation}
  0
  \label{eq:E7}
\end{equation}

\end{document}